% Generated draft: requires mathematical and source review.
\documentclass[11pt]{article}
\usepackage[margin=25mm]{geometry}
\usepackage{fontspec}
\setmainfont{Latin Modern Roman}
\usepackage{amsmath,amssymb,mathtools,unicode-math}
\setmathfont{Latin Modern Math}
\usepackage{enumitem,needspace,xurl,hyperref}
\setlength{\parindent}{0pt}
\setlength{\parskip}{5pt}
\setlength{\emergencystretch}{3em}
\newcommand{\N}{\mathbb{N}}
\newcommand{\Z}{\mathbb{Z}}
\newcommand{\Q}{\mathbb{Q}}
\newcommand{\R}{\mathbb{R}}
\newcommand{\C}{\mathbb{C}}
\newcommand{\F}{\mathbb{F}}
\newcommand{\A}{\mathbb{A}}
\newcommand{\PP}{\mathbb{P}}
\newcommand{\E}{\mathbb{E}}
\newcommand{\eps}{\varepsilon}
\newcommand{\dd}{\,\mathrm{d}}
\begin{document}
\section*{UnsolvedMath ID: 20001113}
\textbf{Source problem code:} catalogue entry
\subsection{Definitions and mathematical statement}
The extracted exponent in the question is corrupted, and a page header is appended. The research aid recovers the intended target as follows: find an explicit family \(f_p:\F_p\times\F_p\to\F_p\), for prime \(p\), and a fixed \(\varepsilon>0\) such that every pair of sources of size on the scale \(p^{1/2}\) has
\[
 |f_p(A\times B)|\geq p^{1/2+\varepsilon}
\]
for sufficiently large \(p\). The exact meaning and constants of the source notation \(|A|,|B|\sim p^{1/2}\) are not given in the extracted text. An explicit comparable-size formulation is to fix \(0<c\leq C<\infty\) and require the assertion for all \(A,B\subseteq\F_p\) with \(c\sqrt p\leq|A|,|B|\leq C\sqrt p\), with the threshold for \(p\) permitted to depend on \(c,C,\varepsilon\).

The supplied aid reports that the fixed polynomial
\[
 f_p(x,y)=x^2+xy=x(x+y)
\]
solves this recovered problem and is Bourgain's published construction, not a new solution. Its quantitative reported input is a Cartesian-incidence estimate: there are absolute constants \(c_0,c_1>0\) such that for equal-size sources \(|A|=|B|=N\) in the range \(N^3\leq c_0p^2\),
\[
 |\{x^2+xy:x\in A,y\in B\}|\geq c_1N^{5/4}.
\]
The aid attributes the argument to the Stevens--de Zeeuw incidence theorem. Pruning comparable sources to equal cardinality therefore gives
\[
 |f_p(A\times B)|\gg_{c,C}p^{5/8}.
\]
Thus for every fixed \(0<\varepsilon<1/8\), including \(\varepsilon=1/9\), the recovered coefficient-one inequality holds for all sufficiently large primes. The strict loss below \(1/8\) accounts for the unspecified positive implied constant.

The aid further reports a distributional corollary, distinct from the image-size assertion alone. Let \(X,Y\) be independent and uniform on equal \(N\)-element sets in the same incidence range, and put \(Z=X(X+Y)\in\F_p\). For a distribution \(Q\) on \(\F_p\), define
\[
 H_\infty(Q)=-\log_2\max_zQ(z),\qquad
 d_{\mathrm{TV}}(\mathcal L(Z),Q)=\tfrac12\sum_{z\in\F_p}|\Pr(Z=z)-Q(z)|.
\]
It claims an absolute argument constant \(C_{\mathrm{inc}}>0\) such that, for every \(0<\tau<1\), some \(Q\) satisfies
\[
 d_{\mathrm{TV}}(\mathcal L(Z),Q)\leq\tau,\qquad
 H_\infty(Q)\geq\tfrac54\log_2N-
 \log_2\!\left(\frac{C_{\mathrm{inc}}}{\tau(1-\tau)}\right).
\]
The underlying collision estimate and the constant are not supplied in full, so this remains a reported candidate corollary requiring verification.
\subsection{Short English statement}
The supplied aid identifies Bourgain's polynomial \(x^2+xy\) as an explicit two-source expander. Its reported \(N^{5/4}\) image bound gives exponent \(5/8\) for square-root-size sources, and hence any fixed gain below \(1/8\) in the recovered question.
\subsection{Sources}
\begin{itemize}
\item Local registry: permanent UnsolvedMath ID \(20001113\), exact public code \texttt{catalogue entry}; \texttt{aim-combinatorics-notes.json}, \texttt{notes[237]}.
\item AIM workshop ``Recent trends in additive combinatorics'', Problem 3.1, attributed to J. Bourgain: \url{https://aimath.org/WWN/additivecomb/additivecomb.pdf}. The extracted tail ``8 COLLECTED BY ERNIE CROOT AND VSEVOLOD F. LEV'' is a page/header artifact, not part of the inequality. The literature field is empty.
\item The supplied statement-verification field reports that the official PDF was checked directly. This is a claim of that field; this draft has not accessed the PDF.
\item Research attempt 1: \texttt{full\_solution}, status at run \texttt{solved}, rigor label \texttt{proved}, novelty confidence low. It attributes the construction to Bourgain's published work and its incidence argument to Stevens--de Zeeuw, and describes the min-entropy statement as a candidate corollary. No full Bourgain or Stevens--de Zeeuw bibliographic entries are supplied. The current registry status is solved and classification \texttt{SOLVED-IN-LITERATURE}.
\end{itemize}
\end{document}
